\par\SourceLeaf{433}%
\noindent\hspace*{1.2em}Um dessen Grenze zu bestimmen, projiciren wir zuerst beide auf die Ebene der $x,y$; die Grössen der Projectionen werden $Zd\sigma$, $Zd\Sigma$ sein: das Zeichen von $Z$ wird den gleichen oder entgegengesetzten Sinn, in welchem die Grenzen die Flächen und ihre Projectionen umgeben, ausdrücken. Man stelle sich die Figur als ein Dreieck vor, die Projection auf die Ebene der $x,y$ habe die Coordinaten
\[
x,y;\qquad x+dx,\ y+dy;\qquad x+\delta x,\ y+\delta y;
\]
und so wird sein doppelter Inhalt
\[
2Zd\sigma=dx.\delta y-dy.\delta x
\]
sein. Der Projection des correspondirenden Elements auf der Kugelfläche werden die Coordinaten entsprechen
\[
\begin{array}{cc}
X,&Y,\\[4pt]
X+\dfrac{\partial X}{\partial x}.dx+\dfrac{\partial X}{\partial y}.dy,&Y+\dfrac{\partial Y}{\partial x}.dx+\dfrac{\partial Y}{\partial y}.dy,\\[6pt]
X+\dfrac{\partial X}{\partial x}.\delta x+\dfrac{\partial X}{\partial y}.\delta y,&Y+\dfrac{\partial Y}{\partial x}.\delta x+\dfrac{\partial Y}{\partial y}.\delta y.
\end{array}
\]
Daraus findet sich die doppelte Area des Elements
\[
\begin{aligned}
2Zd\Sigma={}&\left(\frac{\partial X}{\partial x}.dx+\frac{\partial X}{\partial y}.dy\right)\left(\frac{\partial Y}{\partial x}.\delta x+\frac{\partial Y}{\partial y}.\delta y\right)\\
&-\left(\frac{\partial X}{\partial x}.\delta x+\frac{\partial X}{\partial y}.\delta y\right)\left(\frac{\partial Y}{\partial x}.dx+\frac{\partial Y}{\partial y}.dy\right)\\
={}&\left(\frac{\partial X}{\partial x}.\frac{\partial Y}{\partial y}-\frac{\partial X}{\partial y}.\frac{\partial Y}{\partial x}\right)(dx.\delta y-dy.\delta x).
\end{aligned}
\]
Das Krümmungsmass ist also
\[
=\frac{\partial X}{\partial x}\frac{\partial Y}{\partial y}-\frac{\partial X}{\partial y}\frac{\partial Y}{\partial x}=\omega.
\]
Nun hat man, da
\[
\begin{gathered}
X=-tZ,\qquad Y=-uZ,\\
(1+tt+uu)ZZ=1:
\end{gathered}
\]
\[
\begin{aligned}
dX&=-Z^{3}(1+uu)dt+Z^{3}tu.du\\
dY&=+Z^{3}tu.dt-Z^{3}(1+tt)du,
\end{aligned}
\]
